\todo{Lecture 19}
\section{Capacite et conducteur}
\subsection{Definition de la capacite}

On a deux conducteurs initialement non-charges ($\phi _{A}=\phi _{B}=0$), une charge $-q$ enlevée de $A$ et placée sur $B$. Ceci induit une difference de potentiel
$$
\phi(A)-\phi(B)=-\int_{b}^{a} \boldsymbol{E} \, d\boldsymbol{\ell}>0.
$$
On observe aussi que $\boldsymbol{E}(\boldsymbol{r})$ est proportionnel a $q$ et donc a $\phi(A)-\phi(B)$ aussi. En particulier
$$
\frac{q}{\phi(A)-\phi(B)}=C.
$$
Souvent, on écrit $\phi(A)-\phi(B)=U$ donc 
$$
C=\frac{q}{U}
$$
ou la constante $C$ est la capacite entre deux conducteurs. On mesure $C$ en $\mathrm{CV}^{-1}=\mathrm{C}^{2}\mathrm{s}^{2}\mathrm{kg}^{-1}\mathrm{m}^{-2}=\mathrm{A}^{2}\mathrm{s}^{4}\mathrm{kg}^{-1}\mathrm{m}^{-2}=\mathrm{F}$ farads.

Si les conducteurs sont entoures du vide, $C$ ne depend que de la géométrie des deux conducteurs.

\subsection{Le condensateur}
\begin{definition}
Un condensateur est un système de deux conducteurs avec charge $\pm q$.
\end{definition}
Normalement, la géométrie est choisie pour maximiser le rapport capacite et volume.
\begin{example}[Condensateur cylindrique]
\begin{figure}[h!]
\centering
\includegraphics[scale=0.4]{condensateur cyl.png}
\end{figure}
On suppose $\ell \gg R_{2}-R_{1}$, donc les effets de bord sont négligeables. On a par Gauss
$$
\underbrace{ \iint _{S}\boldsymbol{E}\, d\boldsymbol{S}}_{ \approx 2\pi r\ell E(r) }=\frac{q}{\varepsilon_{0}}\Rightarrow E(r)\approx \frac{q}{\varepsilon_{0}2\pi r\ell}.
$$
On calcule 
\begin{align*}
\phi(B)-\phi(A) & =\int_{R_{1}}^{R_{2}} \boldsymbol{E} \, d\boldsymbol{\ell} \\
  & =-\int_{R_{1}}^{R_{2}} E(r)\,\boldsymbol{\hat{e}}_{r}\cdot \boldsymbol{\hat{e}}_{r} \, dr \\
 & =-\int_{R_{1}}^{R_{2}} E(r) \, dr \\
  & =-\frac{q}{2\pi\varepsilon_{0}\ell}\int_{R_{1}}^{R_{2}} \frac{1}{r} \, dr  \\
 & =-\frac{q}{2\pi\varepsilon_{0}\ell}\log\left( \frac{R_{2}}{R_{1}} \right) \\
 \Rightarrow C & =\frac{q}{\phi(A)-\phi(B)}=\frac{2\pi\varepsilon_{0}\ell}{\log(R_{2}/R_{1})}
\end{align*}
\end{example}
\begin{example}[Condensateur plan]
\begin{figure}[h!]
\centering
\includegraphics[scale=0.6]{condensateur plan.png}
\end{figure}
On note l'aire de la face interieure d'une plaque par $S_{c}=ab$. On suppose $d\ll a$ et $d\ll b$, donc les effets de bord sont negligeables. Par Gauss
$$
\underbrace{ \iint _{S}\boldsymbol{E}~d\boldsymbol{S} }_{ \approx ES_{c} }=\frac{q}{\varepsilon}\Rightarrow E= \frac{q}{\varepsilon_{0}S_{c}}.
$$
On calcule
\begin{align*}
\phi(B)-\phi (A) & =-\int _{A}^{B}\boldsymbol{E}\,d\boldsymbol{\ell} \\
 & =-\int_{A}^{B}E(x)\,\boldsymbol{\hat{e}}_{x}\cdot \boldsymbol{\hat{e}}_{x}  \, dx \\
 & =-\int_{0}^{b} E \, dx =-Ed \\
 & =-\frac{qd}{\varepsilon_{0}S_{c}} \\
 \Rightarrow C & =\frac{q}{\phi(A)-\phi(B)}= \frac{\varepsilon_{0}S_{c}}{d} .
\end{align*}
\end{example}

\subsection{Energie stockée dans un conducteur et densité d'énergie électrique}

Soient deux condensateurs arbitraires. On définit $U=\phi(A)-\phi(B)$. Le travail pour déplacer une charge $-dq$ du conducteur $A$ au conducteur $B$ est 
\begin{align*}
dW & =(\phi(B)-\phi(A))(-dq) \\
 & =(\phi(A)-\phi(B))dq \\
 & =Udq=\frac{q}{C}dq
\end{align*}

Alors, le travail pour charger le conducteur a la charge $q$ est
\begin{align*}
W & =\int_{0}^{q} dW \\
 & =\int_{0}^{q} \frac{1}{C}q' \, dq' \\
 &  =\frac{q^{2}}{2C}=\frac{1}{2}CU^{2} 
\end{align*}`'
Pour un condensateur plan, on a 
$$
W=\frac{1}{2}CU^{2}=\frac{1}{2} \frac{\varepsilon_{0}S_{c}}{d}(Ed)^{2}=\frac{1}{2}\varepsilon_{0}E^{2}S_{c}d
$$
ou $S_{c}d$ est le volume entre les deux plaques (la region ou $\boldsymbol{E}\neq \boldsymbol{0}$).

Ce travail est stockée dans le champ électrique $\boldsymbol{E}$. On définit
$$
e_{E}=\frac{1}{2}\varepsilon_{0}E^{2}
$$
la densité d'énergie électrique (en absence d'un diélectrique); valable en générale pour chaque champ $\boldsymbol{E}$, pas seulement dans un condensateur plan.  

\subsection{Capacite avec un diélectrique}
\begin{definition}
Un diélectrique est un isolant, il ne contient pas de porteurs de charge libre.
\end{definition}
\begin{definition}
Le moment dipolaire $\boldsymbol{p}$ pour un systeme de charges globalement neutre (ici une molecule). Dans le cas de charges ponctuelles on le definit comme
$$
\boldsymbol{p}=\sum _{i}\boldsymbol{r}_{i}q_{i}.
$$
Comme le systeme est globalement neutre, $\sum _{i}q_{i}=0$.
\end{definition}
\begin{example}
Soit $\boldsymbol{p}=-\boldsymbol{r}_{1}q+\boldsymbol{r}_{2}q=q(\boldsymbol{r}_{2}-\boldsymbol{r}_{1})$. Alors $\boldsymbol{p}$ est nul si $\boldsymbol{r}_{2}=\boldsymbol{r}_{1}$, sinon il pointe de la charge negative vers la charge positive.
\end{example}
\paragraph{Properties}
\begin{enumerate}
\item $\boldsymbol{p}$ ne depend pas de l'origine du système de coordonnées. 
\item Dans un champ $\boldsymbol{E}$ uniforme, la force sur le système de charges est nulle. Car $\boldsymbol{F}_{\mathrm{res}}=\sum _{i}q_{i}\boldsymbol{E}=\boldsymbol{E}\sum _{i}q_{i}=0$.
\item En générale, le moment de force n'est pas nul. Soit $\boldsymbol{E}$ uniforme, 
$$
\boldsymbol{M}_{\mathrm{res}}=\sum _{i}\boldsymbol{r}_{i}\times \boldsymbol{F}_{i}=\sum _{i}\boldsymbol{r}_{i}\times \boldsymbol{E}q_{i}=\left( \sum _{i}\boldsymbol{r}_{i}q_{i} \right)\times \boldsymbol{E}=\boldsymbol{p}\times \boldsymbol{E}
$$
ce moment de force tend a orienter $\boldsymbol{p}$ le long de $\boldsymbol{E}$.
\end{enumerate}



